\documentclass[12pt]{article}
\usepackage[utf8]{inputenc}
\usepackage[T1]{fontenc}
\usepackage{hyperref}
\title{Walloon Language Resources for K-12 Curriculum Development}
\author{Compiled by Claude Sonnet 5.5 (claude-sonnet-5-5) for Andr\'as Kornai}
\date{}
\begin{document}
\maketitle

\section{Introduction}
Walloon (ISO 639-3 wln) is a Romance language of southern Belgium. This search was shallow and found only a few reliable sources. Several generic ``skill guide'' sites turned up in results and are not reliable, so they are not cited.

\section{Primary Educational Resources}
The Common Written Walloon (Rifondou Walon) site is in French and lists a resource on teaching Walloon to young schoolchildren who are already learning French. I did not open it. \url{https://rifondou.walon.org/index-engl.html}

\section{Linguistic and Reference Materials}
\begin{itemize}
\item Rifondou Walon (Common Written Walloon). The search summary says it lists an explanatory online Walloon dictionary, an English--Walloon and Walloon--English lexicon, an online grammar, a small grammar written in Walloon, and conjugation tables. \url{https://rifondou.walon.org/index-engl.html}
\item Universit\'e catholique de Louvain course listing on Walloon dialect languages and literature (2005, taught by Jean Germain, offered every two years, so it may not be current): \url{https://sites.uclouvain.be/etudes/2005/cours/en/rom2471a.pdf}
\end{itemize}

\section{Religious and Community Resources}
Nothing specific was found. The Walloon-language forum below is a community discussion space: \url{https://berdelaedje.walon.org/viewtopic.php?p=1945}. A 2009 learner post on it says few learning materials exist and that older textbooks are out of print.

\section{Digital and Audio Resources}
Only the online dictionary, lexicon and grammar on the Rifondou site were seen. No audio or video resources were found.

\section{Curriculum Development Context}
Walloon has many speakers but few learning materials in this search. The Rifondou site is the most promising starting point for a curriculum. Its contents should be checked directly, because I only saw a summary.
\end{document}